\section{Introduction}
%\label{section: introduction}

Let $\bigl(\mathbb{S}^2,g_0\bigr)$ be the $2$-dimensional
sphere with the standard smooth metric $g_0$.
In this paper, we consider the following classical problem in conformal geometry:
Given $t_1,\dots,t_N\in \mathbb{S}^2$ and a set
of positive real numbers $\theta_{t_1},\dots,\theta_{t_N}$, is there
 a metric of constant curvature $1$ conformal to $g_0$
such that the
metric is smooth on $\mathbb{S}^2\setminus\{t_1,\dots,t_N\}$ and has conical
singularities at $t_j$ with the angle $2\pi\theta_{t_j}$? This problem has been widely studied in the literature, and in particular, it
was solved in \cite{Eremenko,Coaxial,EGT} in the cases where there are at
most $2$ non-integer angles among
$\{\theta_{t_1},\dots,\theta_{t_N}\}$ \big(We simply say that the angle
is non-integer if $\theta_j\notin\mathbb{Z}$\big). In this paper, we
consider the case that there are at least $3$ non-integer angles among
$\{\theta_{t_1},\dots,\theta_{t_N}\}$. Without loss of generality, we
may assume that three singularities with non-integer angles are $0$,
$1$, and $\infty$.
Then it is well known that this problem is equivalent to solving
 the following curvature
equation
\begin{gather}
 \Delta u+{\rm e}^u=4\pi\Biggl(\alpha_0\delta_{0}+\alpha_1\delta_{1}
 +\sum_{j=1}^{n+m}\alpha_{t_j}\delta_{t_j}\Biggr)
 \qquad \text{in }\mathbb R^2\cong\C,\nonumber \\
 u(x)=-2(2+\alpha_\infty)\ln|x|+O(1) \qquad \text{as }|x|\to\infty, \label{00eq: DE on C}
 \end{gather}
where $\delta_p$ is the Dirac measure at $p$, $n\geq 0$, $m\geq 1$, $\alpha_0, \alpha_1, \alpha_\infty,
\alpha_{t_j}\in\R_{>-1}\setminus\mathbb{Z}$ for all $1\leq j\leq n$
and $\alpha_{t_j}\in\mathbb{N}$ for all $n+1\leq j\leq
m$. Geometrically, a solution $u(x)$ of (\ref{00eq: DE on C}) leads to
a conformal metric ${\rm d}s^2=\frac12{\rm e}^u|{\rm d}x|^2$ of
curvature $1$ on $\mathbb{S}^2$ with angles $2\pi\theta_{p}$ at conical
singularities $p\in\{0,1,\infty, t_j\}$, where the angles are given by
\begin{equation*} %\label{00eq: theta alpha appendix}
 \theta_{p}=\alpha_p+1, \qquad
 p\in \{0,1,\infty, t_1,\dots,t_{n+m}\},
\end{equation*}
so
\begin{gather*} %\label{00eq: theta alpha a}
 \theta_{p}\in \mathbb{N}_{\geq 2}, \quad
 p\in \{t_{n+1},\dots,t_{n+m}\},\qquad
 \theta_p\in \R_{>0}\setminus\mathbb{Z},\quad p\in\{0,1,\infty,t_1,\dots,t_n\}.
\end{gather*}
It is known that if~\eqref{00eq: DE on C} has a solution $u(x)$, then
there is a developing map $h(x)$ such that $h(x)$ is a multi-valued meromorphic
function on $\C$ satisfying
\begin{equation*}%\label{eqLivoulle}
 u(x)=\ln\frac{8|h'(x)|^2}{(1+|h(x)|^2)^2},\qquad x\in
 \dot{\C}:=\mathbb{C}\setminus\{0,1,t_1,\dots,t_{m+n}\}.
\end{equation*}
This is known as the Liouville formula and this $h(x)$ is multi-valued
and has its projective monodromy group contained in $\PSU(2)$. More
precisely, given any $\ell\in \pi_1\bigl(\dot{\C}, q_0\bigr)$ where
$q_0\in\dot{\C}$ is base point, the analytic continuation of $h(x)$,
denoted by $\ell^*h(x)$, is also a developing map of~$u(x)$ and so
$\ell^*h(x)=\frac{ah(x)+b}{ch(x)+d}$ for some projective monodromy
matrix $\SM{a}{b}{c}{d}\in\PSU(2)$.

The curvature equation~\eqref{00eq: DE on C} has been extensively
studied in the past several decades. For the recent development of
this subject and its application, we refer
\cite{Chai-Lin-Wang,Chen-Lin-weight2,
 Chen-Lin-rectangular,Chen-Lin-Green-function,Chen-Lin-Yang,
Eremenko,Coaxial,EGT,Quadrilateral,Guo-Lin-Yang,
 Lin-Wang-bubbling,Luo-Tian,Mondello-Panov,Mondello-Panov2}
to the interested reader. So far, the existence of solutions for
\eqref{00eq: DE on C} with general singular sources is challenging
and seems to be out of reach.
Mondello and Panov \cite{Mondello-Panov} studied a reduced
problem: \emph{To describe possible angles $($i.e., to describe the data
 $\{\theta_0,
 \theta_1,\theta_\infty,\theta_{t_1},\dots,\theta_{t_{n+m}}\})$ such
 that~\eqref{00eq: DE on C} has solutions for some singular set
 $\{t_1,\dots,t_{n+m}\}$}.
They introduced the measure
$d_1\bigl(\Z_o^{n+m+3},{\boldsymbol\theta-\boldsymbol
 1}\bigr)=d_1\bigl(\Z_o^{n+m+3},{\boldsymbol\alpha}\bigr)$ on
\[
 \boldsymbol \theta=(\theta_0,\theta_1,\theta_\infty,\theta_{t_1},
 \dots,\theta_{t_{n+m}}),
\]
i.e.,
\[
 \boldsymbol\theta-\boldsymbol 1={\boldsymbol\alpha}
 =(\alpha_0,\alpha_1,\alpha_\infty,\alpha_{t_1},\dots,
 \alpha_{t_{n+m}}),
\]
where $\Z_o^{n+m+3}$ is the subset of $\Z^{n+m+3}$ consisting of
vectors with odd sums of coordinates, and~$d_1$~is the
$\ell_1$-distance. Then they proved the following interesting result.

\begin{Theorem}[{\cite{Mondello-Panov}}]\label{thm-MP}
 Suppose that the angles $\theta_0$, $\theta_1$, $\theta_\infty$ and
 $\theta_{t_i}$, $1\le i\le n+m$, are fixed. Then the following
 hold.
\begin{enumerate}\itemsep=0pt
\item[$1.$]
If~\eqref{00eq: DE on C} has a solution $u(x)$ for some singular set $\{t_1,\dots,t_{n+m}\}$, then
\[d_1\bigl(\Z_o^{n+m+3},{\boldsymbol\theta-\boldsymbol 1}\bigr)\geq 1.\]
Furthermore, if
 \begin{equation}\label{00e0q-5}
d_1\bigl(\Z_o^{n+m+3},{\boldsymbol\theta-\boldsymbol 1}\bigr)=1,
\end{equation}
then there is a developing map $h(x)$ such that the projective
monodromy group of $h(x)$ is contained in the unit circle, i.e., any
projective monodromy matrix of $h(x)$ is diagonal.

\item[$2.$]
Conversely, if $d_1\bigl(\Z_o^{n+m+3},{\boldsymbol\theta-\boldsymbol 1}\bigr)>1$, then there exists some singular set $\{t_1,\dots, t_{n+m}\}$ such that~\eqref{00eq: DE on C} has solutions.
\end{enumerate}
\end{Theorem}

\begin{Definition}
In \cite{Mondello-Panov}, the corresponding metric
$\frac12{\rm e}^{u(x)}|{\rm d}x|^2$ (if exists) is called \emph{co-axial} if there
is a developing map $h(x)$ such that the projective monodromy group of
$h(x)$ is contained in the unit circle. In this paper, we also call
this solution $u(x)$ \emph{co-axial} for convenience.
\end{Definition}

For the case~\eqref{00e0q-5}, Theorem~\ref{thm-MP} shows that solutions must be
co-axial if exist. However, the existence of co-axial metrics does not
necessarily imply~\eqref{00e0q-5}. Thus, there arises naturally the
following question: Are there sufficient conditions on angles that
guarantee the existence of a~singular set $\{t_1,\dots,t_{n+m}\}$ such that
\eqref{00eq: DE on C} has a co-axial solution? This question was
completely solved by Eremenko \cite{Coaxial}.


\begin{Theorem}[\cite{Coaxial}]\label{thm-Eremenko}
If~\eqref{00eq: DE on C} has a co-axial solution $u(x)$ for some singular set $\{t_1,\dots,t_{n+m}\}$, then
 there are $\epsilon_{p}\in\{\pm 1\}$ for $p\in
 \{0,1,\infty,t_1,\dots,t_n\}$ such that
\begin{gather}
k':=\sum_{p\in\{0,1,\infty,t_1,\dots,t_n\}} \epsilon_{p}\theta_{p}\in\mathbb{Z}_{\geq 0},\label{00e0qq-111}
\\
k'':=\sum_{j=n+1}^{n+m} \theta_{t_j}-(n+m+3)-k'+2\in 2\mathbb{Z}_{\geq 0}.\label{00e0qq-112}
\end{gather}
Moreover, if there is $\eta\in\mathbb{R}_{>0}$ such that
$\mathbf{c}=\eta\mathbf{b}=\eta(b_1,\dots,b_q)$ with
$b_j\in\mathbb{N}$, $\gcd(b_1,\dots,b_q)=1$, where
\[
 \boldsymbol c:=(\theta_0,\theta_1,\theta_\infty,\theta_{t_1},
 \dots,\theta_{t_n},\underbrace{1,\dots,1}_{k'+k''}),
\]
then we also have
\begin{equation}\label{k3}
 2\max_{n+1\leq j\leq {n+m}}\theta_{t_j}\leq \sum_{j=1}^q b_j.
\end{equation}



Conversely, if there is no such $\eta$ such that
$\mathbf{c}=\eta\mathbf{b}$, then
\eqref{00e0qq-111} and~\eqref{00e0qq-112} are also sufficient; if there
is such $\eta$ such that $\mathbf{c}=\eta\mathbf{b}$, then
\eqref{00e0qq-111}--\eqref{k3} are also sufficient.
\end{Theorem}

As we have remarked earlier, the existence of co-axial solutions does
not necessarily imply~\eqref{00e0q-5}. Let us recall Eremenko's
example \cite{Coaxial}:
$(\theta_0,\theta_1,\theta_\infty,\theta_{t_1},\theta_{t_2})=(\beta,\beta,2\beta,2\beta,3)$
for some $\beta\in\mathbb{R}_{>0}\setminus\N$, for which it follows
from Theorem~\ref{thm-Eremenko} that co-axial solutions exist but it
does not satisfy~\eqref{00e0q-5}.

Here we have an interesting observation about Mondello--Panov's
condition~\eqref{00e0q-5} and Eremenko's condition~\eqref{k3}.

\begin{Theorem}\label{thm-MP-E}
Suppose that~\eqref{00e0q-5} holds. Then
\eqref{00e0qq-111} and~\eqref{00e0qq-112} imply~\eqref{k3} automatically,
namely~\eqref{00eq: DE on C} has $($co-axial$)$ solutions for some
singular set $\{t_1,\dots,t_{n+m}\}$ if and only if
\eqref{00e0qq-111} and~\eqref{00e0qq-112} hold for some
$\epsilon_p\in\{\pm 1\}$.
\end{Theorem}

Theorem~\ref{thm-MP-E} will be proved in Section~\ref{sec-Eremenko}.
Define
\[
 \triangle_{n+m}:=\bigl\{(t_1,\dots,t_{n+m})\in \mathbb{C}^{n+m}\mid
 t_j=0,1 \text{ or } t_j=t_k\;\text{for some $j\neq k$}\bigr\}.
\]
Throughout the paper, we let $\theta_0$, $\theta_1$,
$\theta_\infty$, and $\theta_j'$, $j=1,\dots,n+m$, be fixed real
numbers such that $\theta_0$, $\theta_1$, $\theta_\infty$, and
$\theta_j'$, $j=1,\dots,n$, are non-integers and $\theta_j'$,
$j=n+1,\dots,n+m$, are integers greater than $1$.
Let
\begin{gather}
 A=A_\theta:=\bigl\{(t_1,\dots,t_{n+m})\in \mathbb{C}^{n+m}\setminus
 \triangle_{n+m}\mid \text{\eqref{00eq: DE on C}
 with}\ \theta_{t_j}=\theta_j',
\nonumber\\ \hphantom{ A=A_\theta:=\{}{}
j=1,\dots,n+m, \
 \text{has co-axial solutions}\bigr\}.\label{eq: set A}
\end{gather}
An $(n+m)$-tuple $(t_1,\dots,t_{n+m})$ in $A$ will be called
admissible for~\eqref{00eq: DE on C}.
By Theorem~\ref{thm-Eremenko}, we know that $A\neq\varnothing$ if and
only if~\eqref{00e0qq-111}--\eqref{k3} hold. In this case, a natural
question that interest us is:
\[\parbox{\dimexpr\linewidth-5em}{\it
 What are the geometric and algebaric properties of the set $A$, such as its dimension and so on?
 }\]
This paper aims to study this problem. Note that~\eqref{00e0qq-112}
implies $m\geq 1$, this is the reason why we assume $m\geq 1$ in this paper.

Let us consider the special case $m=1$ first. Denote
\[
 \Q_{\theta}:=\Q\bigl(\theta_0,\theta_1,\theta_\infty,\theta_1',
 \dots,\theta_n'\bigr)\subsetneqq \R.
\]

\begin{Theorem}\label{00thm-2-2}
Let $m=1$ and~\eqref{00e0qq-111}--\eqref{k3} hold with
$\theta_{t_j}=\theta_j'$, i.e., $A\neq\varnothing$.
Then $A$ is a finite set with $\# A\leq
2^{n+2}\times\bigl(\theta_{n+1}'-1\bigr)!$. Furthermore, for each
$\boldsymbol t=(t_1,\dots,t_{n+1})\in A$, all $t_j$'s are algebraic
over $\Q_{\theta}$ and the field
\[\Q_\theta(A):=\Q_{\theta}\bigl(\bigl\{t_j \mid (t_1,\dots,t_{n+1})\in A\bigr\}\bigr)\]
is a Galois extension of $\Q_{\theta}$ and
\[[\Q_\theta(A) : \Q_\theta]\leq M<\infty,\]
where $M$ is a constant depending on $n$ and $\theta_{n+1}'$.
\end{Theorem}

In the special case when there are only three
non-integer angles, i.e., when $n=0$, we can get a much sharper bound
for the cardinality of $A$. In the following, we will write
$t_1$ and $\theta_1'$ simply by $t$ and $\theta$.

\begin{Theorem} \label{theorem: sharp}
 Assume that $n=0$ and $m=1$. Assume that
 $\epsilon_0,\epsilon_1\in\{\pm1\}$ are signs such that
 $k:=\theta_\infty+\epsilon_0\theta_0+\epsilon_1\theta_1\in\Z$ and
 $\theta\equiv k\mod 2$. Then the cardinality of the set $A$ is
 \[
 \begin{cases}
 0, &\text{if }\theta\le|k|, \\
 \bigl(\theta^2-k^2\bigr)/4, &\text{if }\theta>|k|, \end{cases}
 \]
 counted with multiplicities $($see Section~$\ref{section: proof of
 Theorem 1.4}$ for an explanation of
 multiplicities$)$. Moreover, if ${t\in A}$, then the degree of $t$ over
 $\Q_\theta$ is at most $\bigl(\theta^2-k^2\bigr)/4$.
\end{Theorem}

We remark that the condition $\theta\equiv k\mod 2$ is necessary for
$A\neq\varnothing$, by~\eqref{00e0qq-112}. Interestingly, a key
ingredient in the proof of the theorem is the monodromy of the
classical hypergeometric equations.


